\section{AMC 12 12A 2002}
        \begin{enumerate}
        	\item (\textit{AMC 12 12A 2002}) Compute the sum of all the roots of
$(2x+3)(x-4)+(2x+3)(x-6)=0$
$\mathrm{(A) \ } \frac{7}{2}\qquad \mathrm{(B) \ } 4\qquad \mathrm{(C) \ } 5\qquad \mathrm{(D) \ } 7\qquad \mathrm{(E) \ } 13$

 \textbf{Answer: }A



\textbf{Solution 1}

We expand to get $2x^2-8x+3x-12+2x^2-12x+3x-18=0$ which is $4x^2-14x-30=0$ after combining like terms. Using the quadratic part of \href{https://artofproblemsolving.com/wiki/index.php?title=Vieta\%27s\_Formulas}{Vieta's Formulas}, we find the sum of the roots is $\frac{14}4 = \boxed{\textbf{(A) }7/2}$.

 



\textbf{Solution 2}

We see a $(2x+3)$ term in both addends so we can factor the equation to $(2x+3)((x-4)+(x-6))=(2x+3)(2x-10)=0.$ The two roots become obvious: they are $-\frac{3}{2}$ and $\frac{10}{2}.$ Adding these two gives $\boxed{\textbf{(A) }7/2}.$

 -starwarmonkey

 


	\item (\textit{AMC 12 12A 2002}) Cindy was asked by her teacher to subtract 3 from a certain number and then divide the result by 9. Instead, she subtracted 9 and then divided the result by 3, giving an answer of 43. What would her answer have been had she worked the problem correctly?


 $\mathrm{(A) \ } 15\qquad \mathrm{(B) \ } 34\qquad \mathrm{(C) \ } 43\qquad \mathrm{(D) \ } 51\qquad \mathrm{(E) \ } 138$

 \textbf{Answer: }A



\textbf{Solution}

We work backwards; the number that Cindy started with is $3(43)+9=138$. Now, the correct result is $\frac{138-3}{9}=\frac{135}{9}=15$. Our answer is $\boxed{\textbf{(A) }15}$.

 


 



\textbf{Solution 2}

Let the number be $x$. We transform the problem into an equation: $\frac{x-9}{3}=43$. Solve for $x$ gives us $x=138$. Therefore, the correct result is $\frac{138-3}{9}=\frac{135}{9}=\boxed{\textbf{(A) }15}$.

 


	\item (\textit{AMC 12 12A 2002}) According to the standard convention for exponentiation, 

 \[2^{2^{2^{2}}} = 2^{\left(2^{\left(2^2\right)}\right)} = 2^{16} = 65536.\]


 If the order in which the exponentiations are performed is changed, how many other values are possible?


 $\mathrm{(A) \ } 0\qquad \mathrm{(B) \ } 1\qquad \mathrm{(C) \ } 2\qquad \mathrm{(D) \ } 3\qquad \mathrm{(E) \ } 4$

 \textbf{Answer: }B



\textbf{Solution 1}

The best way to solve this problem is by simple brute force. 

 It is convenient to drop the usual way how exponentiation is denoted, and to write the formula as $2\uparrow 2\uparrow 2\uparrow 2$, where $\uparrow$ denotes exponentiation. We are now examining all ways to add parentheses to this expression. There are 5 ways to do so:

 
            \begin{enumerate}
        \item $2\uparrow (2\uparrow (2\uparrow 2))$
        \item $2\uparrow ((2\uparrow 2)\uparrow 2)$
        \item $((2\uparrow 2)\uparrow 2)\uparrow 2$
        \item $(2\uparrow (2\uparrow 2))\uparrow 2$
        \item $(2\uparrow 2)\uparrow (2\uparrow 2)$
            \end{enumerate}
            We can note that $2\uparrow (2\uparrow 2) = (2\uparrow 2)\uparrow 2 =16$. Therefore options 1 and 2 are equal, and options 3 and 4 are equal.Option 1 is the one given in the problem statement. Thus we only need to evaluate options 3 and 5.

 $((2\uparrow 2)\uparrow 2)\uparrow 2 = 16\uparrow 2 = 256$
$(2\uparrow 2)\uparrow (2\uparrow 2) = 4 \uparrow 4 = 256$

 Thus the only other result is $256$, and our answer is $\boxed{\textbf{(B) } 1}$.

 



\textbf{Solution 2 (Recursive Method)}

We will proceed by recursion using the same notation as Solution 1. Note that we can either have $(2 \uparrow 2 \uparrow 2) \uparrow 2$ or $2 \uparrow (2 \uparrow 2 \uparrow 2)$.

 However, the expression in the parentheses can either be $(2 \uparrow 2) \uparrow 2$ or $2 \uparrow (2 \uparrow 2)$, which correspond to the same value of $16$. Therefore, we can either have $16^2$ or $2^{16}$, so our answer is $\boxed{\textbf{(B) }1}$.

 


	\item (\textit{AMC 12 12A 2002}) Find the degree measure of an angle whose complement is 25\% of its supplement.


 $\mathrm{(A) \ 48 } \qquad \mathrm{(B) \ 60 } \qquad \mathrm{(C) \ 75 } \qquad \mathrm{(D) \ 120 } \qquad \mathrm{(E) \ 150 }$

 \textbf{Answer: }B



\textbf{Solution 1}

We can create an equation for the question, $4(90-x)=(180-x)$
$360-4x=180-x$
$3x=180$

 After simplifying, we get $x=60 \Rightarrow \mathrm {(B)}$




\textbf{Solution 2}

Given that the complementary angle is $\frac{1}{4}$ of the supplementary angle. Subtracting the complementary angle from the supplementary angle, we have $90^{\circ}$ as $\frac{3}{4}$ of the supplementary angle.

 Thus the degree measure of the supplementary angle is $120^{\circ}$, and the degree measure of the desired angle is $180^{\circ} - 120^{\circ} = 60^{\circ}$. $\mathrm {(B)}$



	\item (\textit{AMC 12 12A 2002}) Each of the small circles in the figure has radius one. The innermost circle is tangent to the six circles that surround it, and each of those circles is tangent to the large circle and to its small-circle neighbors. Find the area of the shaded region.


 
                    \begin{figure}[H]
                        \centering
                        \includegraphics[width=0.5\textwidth]{Problem/AMC_12_12A_2002_5_Problem_Figure_1.png}
                    \end{figure}
                    
$\text{(A)}\ \pi \qquad \text{(B)}\ 1.5\pi \qquad \text{(C)}\ 2\pi \qquad \text{(D)}\ 3\pi \qquad \text{(E)}\ 3.5\pi$

 \textbf{Answer: }C



\textbf{Solution}

The outer circle has radius $1+1+1=3$, and thus area $9\pi$. The little circles have area $\pi$ each; since there are 7, their total area is $7\pi$. Thus, our answer is $9\pi-7\pi=\boxed{2\pi\Rightarrow \textbf{(C)}}$.

 


	\item (\textit{AMC 12 12A 2002}) For how many positive integers $m$ does there exist at least one positive integer $n$ such that $m \cdot n \le m + n$?


 $\mathrm{(A) \ } 4\qquad \mathrm{(B) \ } 6\qquad \mathrm{(C) \ } 9\qquad \mathrm{(D) \ } 12\qquad \mathrm{(E) \ }$  infinitely many


 
 \textbf{Answer: }E



\textbf{Solution 1}

For any $m$ we can pick $n=1$, we get $m \cdot 1 \le m + 1$,therefore the answer is $\boxed{\textbf{(E) } \text{infinitely many}}$.

 



\textbf{Solution 2}

Another solution, slightly similar to this first one would be using \href{https://artofproblemsolving.com/wiki/index.php?title=Simon\%27s\_Favorite\_Factoring\_Trick}{Simon's Favorite Factoring Trick}.

 $(m-1)(n-1) \leq 1$

 Let $n=1$, then

 $0 \leq 1$

 This means that there are infinitely many numbers $m$ that can satisfy the inequality. So the answer is $\boxed{\textbf{(E) } \text{infinitely many}}$.

 



\textbf{Solution 3}

If we subtract $n$ from both sides of the equation, we get $m \cdot n - n \le m$. Factor the left side to get $(m - 1)(n) \le m$. Divide both sides by $(m-1)$ and we get $n \le \frac {m}{m-1}$. The fraction $\frac {m}{m-1} > 1$ if $m > 1$. There is an infinite amount of integers greater than $1$, therefore the answer is $\boxed{\textbf{(E) } \text{infinitely many}}$.

 


	\item (\textit{AMC 12 12A 2002}) A $45^\circ$ arc of circle A is equal in length to a $30^\circ$ arc of circle B. What is the ratio of circle A's area and circle B's area?


 $\text{(A)}\ 4/9 \qquad \text{(B)}\ 2/3 \qquad \text{(C)}\ 5/6 \qquad \text{(D)}\ 3/2 \qquad \text{(E)}\ 9/4$

 \textbf{Answer: }A



\textbf{Solution 1}

Let $r_1$ and $r_2$ be the radii of circles $A$ and$B$, respectively. 

 It is well known that in a circle with radius $r$, a subtended arc opposite an angle of $\theta$ degrees has length $\frac{\theta}{360} \cdot 2\pi r$. 

 Using that here, the arc of circle A has length $\frac{45}{360}\cdot2\pi{r_1}=\frac{r_1\pi}{4}$. The arc of circle B has length $\frac{30}{360} \cdot 2\pi{r_2}=\frac{r_2\pi}{6}$. We know that they are equal, so $\frac{r_1\pi}{4}=\frac{r_2\pi}{6}$, so we multiply through and simplify to get $\frac{r_1}{r_2}=\frac{2}{3}$. As all circles are similar to one another, the ratio of the areas is just the square of the ratios of the radii, so our answer is $\boxed{\textbf{(A) } 4/9}$.

 



\textbf{Solution 2}

The arc of circle $A$ is $\frac{45}{30}=\frac{3}{2}$ that of circle $B$.

 The circumference of circle $A$ is $\frac{2}{3}$ that of circle $B$ ($B$ is the larger circle).

 The radius of circle $A$ is $\frac{2}{3}$ that of circle $B$.

 The area of circle $A$ is ${\left(\frac{2}{3}\right)}^2=\boxed{\textbf{(A) } 4/9}$ that of circle $B$.

 



\textbf{Solution 3}

We can use a very cheap way to find the answer based on the choices. Firstly, if a higher angle measure equals a lower angle measure, that means that the ratio of areas must be lower, cutting D and E. Then we realize that the ratio has to be with square numbers, as there are no square roots. Therefore, it is $(A)$

 -dragoon

 



\textbf{Solution 4}


            \begin{itemize}
        \item Written by a 6th grader, please don't mind*
            \end{itemize}
            Set Circle A's arc to 45 degrees. In a full circle, there is 360 degrees, so 360/45 = 8. If we set the arc to be pi degrees, then in total there will be 8 pi in the circumference of this circle, yielding 4 as the radius and 16pi as the area. 

 Since the arc of 30 degrees in Circle B is also pi, and there are 12 30 degree arcs in a circle, there will be 12pi in the circumference of Circle B, yielding 6 as the radius and 36pi as the area. 

 Now, we put 16pi/36pi as our ratio, yielding 16/36=4/9. Therefore, it is $(A)$4/9. 

 



\textbf{Solution 5}

WLOG, let the radius of circle $B$ be $1$, and the radius of circle $A$ be $r$. Arcs of circle $A$ and $B$ are $\frac{45}{180}\pi r$ and $\frac{30}{180}\pi$, respectively. Since arcs are equal length, $\frac{45}{180}\pi r = \frac{30}{180}\pi$ so $r=\frac{2}{3}$. The desired ratio is $\pi r^2 : \pi 1^2$. Substituting, we get $\pi\left(\frac{2}{3}\right)^2 : \pi = \frac{4}{9} : 1 =\boxed{\textbf{(A) } 4/9}$.

 


	\item (\textit{AMC 12 12A 2002}) Betsy designed a flag using blue triangles, small white squares, and a red center square, as shown. Let $B$ be the total area of the blue triangles, $W$ the total area of the white squares, and $R$ the area of the red square. Which of the following is correct?


 
                    \begin{figure}[H]
                        \centering
                        \includegraphics[width=0.5\textwidth]{Problem/AMC_12_12A_2002_8_Problem_Figure_1.png}
                    \end{figure}
                    
$\text{(A)}\ B = W \qquad \text{(B)}\ W = R \qquad \text{(C)}\ B = R \qquad \text{(D)}\ 3B = 2R \qquad \text{(E)}\ 2R = W$

 \textbf{Answer: }A



\textbf{Solution}

The blue that's touching the center red square makes up 8 triangles, or 4 squares. Each of the corners is 1 square and each of the edges is 1, totaling 12 squares. There are 12 white squares, thus we have $\boxed{B=W\Rightarrow \text{(A)}}$.

 



\textbf{Solution 2}

Cut the pattern into 16 square regions and use symmetry.
                    \begin{figure}[H]
                        \centering
                        \includegraphics[width=0.5\textwidth]{Solution/AMC_12_12A_2002_8_Solution_Figure_1.png}
                    \end{figure}
                    

 This clearly shows $\boxed{\textbf{(A) }B=W}$.

 


 



\textbf{Solution 3 (a bit lengthy)}


                    \begin{figure}[H]
                        \centering
                        \includegraphics[width=0.5\textwidth]{Solution/AMC_12_12A_2002_8_Solution_Figure_2.png}
                    \end{figure}
                    

 Suppose one leg of the blue triangle is $x$. Then, the other leg will also be $x$. To prove this, the problem states that the white pieces are squares, which means that all of their sides are congruent. Examine any blue triangle touching a white square using its legs. You can easily see that leg 1 = leg 2. 

 If the value of a leg is $x$, then the hypotenuse becomes $\sqrt(2) x$. Now, count how many hypotenuses make up one side of the red square. You will see that it is $2$. So, the area of the square is $(2 \sqrt(2) x)^2$, or $8x^2$. 

 Next, find the area of a single blue triangle. If its legs are $x$, then the area becomes $(x^2)/2$. Counting the number of triangles, you get 24. Multiply this by the area to get the total area of triangles to make up $12x^2$. 

 If the value of a leg of a triangle is $x$, then so is the value of the side of a white square. This means that the area of a white square is $x^2$. Counting the number of squares, you get 12. This means that the total area of squares is $12x^2$. 

 Because this is the same as the area of triangles, you get: $\boxed{\textbf{(A) }B=W}$.

 -N.

 


	\item (\textit{AMC 12 12A 2002}) Jamal wants to save 30 files onto disks, each with 1.44 MB space. 3 of the files take up 0.8 MB, 12 of the files take up 0.7 MB, and the rest take up 0.4 MB. It is not possible to split a file onto 2 different disks. What is the smallest number of disks needed to store all 30 files?


 $\text{(A)}\ 12 \qquad \text{(B)}\ 13 \qquad \text{(C)}\ 14 \qquad \text{(D)}\ 15 \qquad \text{(E)} 16$

 \textbf{Answer: }B



\textbf{Solution}

A $0.8$ MB file can either be on its own disk, or share it with a $0.4$ MB. Clearly it is better to pick the second possibility. Thus we will have $3$ disks, each with one $0.8$ MB file and one $0.4$ MB file.

 We are left with $12$ files of $0.7$ MB each, and $12$ files of $0.4$ MB each. Their total size is $12\cdot (0.7 + 0.4) = 13.2$ MB. The total capacity of $9$ disks is $9\cdot 1.44 = 12.96$ MB, hence we need at least $10$ more disks. And we can easily verify that $10$ disks are indeed enough: six of them will carry two $0.7$ MB files each, and four will carry three $0.4$ MB files each.

 Thus our answer is $3+10 = \boxed{\textbf{(B) }13 }$.

 



\textbf{Solution 2}

Similarly to Solution 1, we see that there must be $3$ disks to account for the $0.8$ MB file. Secondly, since there are $[30-(3+12)]-3 = 12$ files(for both 0.4 MB and 0.7 MB) left, it is easy to see that the optimal way to place the files would be $3$ files per disks for the 0.4 MB files and hence would require 4 disks. 

 We are left with $12$ files($0.7$ MB), where the optimal number of files per disks is $2$, so the optimal number of disks for this type of file would be $6$ disks. Therefore, the answer is $3+4+6=\boxed{13}$.

 


	\item (\textit{AMC 12 12A 2002}) Sarah places four ounces of coffee into an eight-ounce cup and four ounces of cream into a second cup of the same size. She then pours half the coffee from the first cup to the second and, after stirring thoroughly, pours half the liquid in the second cup back to the first. What fraction of the liquid in the first cup is now cream?


 $\mathrm{(A) \ } \frac{1}{4}\qquad \mathrm{(B) \ } \frac13\qquad \mathrm{(C) \ } \frac38\qquad \mathrm{(D) \ } \frac25\qquad \mathrm{(E) \ } \frac12$

 \textbf{Answer: }D



\textbf{Solution 1}

We will simulate the process in steps.

 In the beginning, we have:

 
            \begin{itemize}
        \item $4$ ounces of coffee in cup $1$
        \item $4$ ounces of cream in cup $2$
            \end{itemize}
            In the first step we pour $4/2=2$ ounces of coffee from cup $1$ to cup $2$, getting:

 
            \begin{itemize}
        \item $2$ ounces of coffee in cup $1$
        \item $2$ ounces of coffee and $4$ ounces of cream in cup $2$
            \end{itemize}
            In the second step we pour $2/2=1$ ounce of coffee and $4/2=2$ ounces of cream from cup $2$ to cup $1$, getting:

 
            \begin{itemize}
        \item $2+1=3$ ounces of coffee and $0+2=2$ ounces of cream in cup $1$
        \item the rest in cup $2$
            \end{itemize}
            Hence at the end we have $3+2=5$ ounces of liquid in cup $1$, and out of these $2$ ounces is cream. Thus the answer is $\boxed{\text{(D) } \frac 25}$.

 



\textbf{Solution 2}

Let's consider this in steps. We have 4 ounces of coffee in the first cup. We hace 4 ounces of cream in the second cup. We take half of the coffee in the first cup(2 ounces), and add it to the second cup, yielding 6 ounces in total in the second cup(in a 1:2 ratio between coffee and cream, respecitvely).We then take half of the second cup and pour it into the first cup. $6/2=3$, so there is now 5 ounces in the first cup, 2 coffee and 3 the mixture.Remember that the mixture is in a 1:2 ratio between coffee and cream. So, coffee has one ounce and cream has 2 ounces. In total, there is 2 ounces in the 5 ounce first cup. Putting 2 over 5, we get the answer. Therefore, the answer is $\boxed{\text{(D) } \frac 25}$.

 


	\item (\textit{AMC 12 12A 2002}) Mr. Earl E. Bird gets up every day at 8:00 AM to go to work. If he drives at an average speed of 40 miles per hour, he will be late by 3 minutes. If he drives at an average speed of 60 miles per hour, he will be early by 3 minutes. How many miles per hour does Mr. Bird need to drive to get to work exactly on time?


 $\text{(A)}\ 45 \qquad \text{(B)}\ 48 \qquad \text{(C)}\ 50 \qquad \text{(D)}\ 55 \qquad \text{(E)} 58$

 \textbf{Answer: }B



\textbf{Solution 1}

Let the time he needs to get there in be $t$ and the distance he travels be $d$. From the given equations, we know that $d=\left(t+\frac{1}{20}\right)40$ and $d=\left(t-\frac{1}{20}\right)60$. Setting the two equal, we have $40t+2=60t-3$ and we find $t=\frac{1}{4}$ of an hour. Substituting t back in, we find $d=12$. From $d=rt$, we find that $r$, our answer, is $\boxed{\textbf{(B) }48 }$.

 



\textbf{Solution 2}

Since either time he arrives at is $3$ minutes from the desired time, the answer is merely the \href{https://artofproblemsolving.com/wiki/index.php?title=Harmonic\_mean}{harmonic mean} of 40 and 60. Substituting $t=\frac ds$ and dividing both sides by $d$, we get $\frac 2s = \frac 1{40} + \frac 1{60}$ hence $s=\boxed{\textbf{(B) }48}$.

 (Note that this approach would work even if the time by which he is late was different from the time by which he is early in the other case - we would simply take a weighted sum in step two, and hence obtain a weighted harmonic mean in step three.)

 



\textbf{Solution 3}

Let x be equal to the total amount of distance he needs to cover. Let y be equal to the amount of time he would travel correctly.Setting up a system of equations, $\frac x{40} -3 = y$ and $\frac x{60} +3 = y$

 Solving, we get x = 720 and y = 15. 

 We divided by y to get the average speeed, $\frac {720}{15} = 48$. Therefore, the answer is $\boxed{\textbf{(B) }48}$.

 



\textbf{Solution 4}

Let $v$ be Mr Bird's speed in miles per hour and $t$ be the desired time in hours. No matter what, the product of Mr Bird's speed and time must always be constant. 

 
                    \begin{figure}[H]
                        \centering
                        \includegraphics[width=0.5\textwidth]{Solution/AMC_12_12A_2002_11_Solution_Figure_1.png}
                    \end{figure}
                    

 From the diagram above, $(60 - v) \left(t - \frac{1}{20} \right) = \frac{v}{20}$ and similarly, $(v - 40)t = \frac{40}{20} = 2$. Expanding, $60t - 3 - vt + \frac{v}{20} = \frac{v}{20}$ and so $60t - 3 - (2 + 40t) = 0, t = \frac{1}{4}$. Hence $v = \frac{2}{1/4} + 40 = 48$.

 


	\item (\textit{AMC 12 12A 2002}) Both roots of the quadratic equation $x^2 - 63x + k = 0$ are prime numbers. The number of possible values of $k$ is 


 $\text{(A)}\ 0 \qquad \text{(B)}\ 1 \qquad \text{(C)}\ 2 \qquad \text{(D)}\ 4 \qquad \text{(E) more than 4}$

 \textbf{Answer: }B



\textbf{Solution 1}

Consider a general quadratic with the coefficient of $x^2$ being $1$ and the roots being $r$ and $s$. It can be factored as $(x-r)(x-s)$ which is just $x^2-(r+s)x+rs$. Thus, the sum of the roots is the negative of the coefficient of $x$ and the product is the constant term. (In general, this leads to \href{https://artofproblemsolving.com/wiki/index.php?title=Vieta\%27s\_Formulas}{Vieta's Formulas}).

 We now have that the sum of the two roots is $63$ while the product is $k$. Since both roots are primes, one must be $2$, otherwise, the sum would be even. That means the other root is $61$ and the product must be $122$. Hence, our answer is $\boxed{\text{(B)}\ 1 }$.

 



\textbf{Solution 2}

Let the quadratic be factored as $(x-r)(x-s)$. Then $-r-s=-63$ and $rs=k$. We know that both $r$ and $s$ are prime, and all primes but $2$ are odd. The sum of two odd numbers is also always even, but in this case the sum of two primes is odd, which means that one of the primes is $2$. This leaves only one possible value for the other root and one possible value for their product $k$.

 



\textbf{Solution 3}

By Vieta's you have 
 \[a+b=63\]
 \[ab=k\]Since we know odd + even = odd, we must have either $a$ or $b$ equal to $2$ and the other equal to $63-2=61.$ Both of these are prime so they satisfy the restraints. Thus there is $\boxed{(B) 1}$ solution.

 


	\item (\textit{AMC 12 12A 2002}) Two different positive numbers $a$ and $b$ each differ from their reciprocals by $1$. What is $a+b$?


 $\text{(A) }1 \qquad \text{(B) }2 \qquad \text{(C) }\sqrt 5 \qquad \text{(D) }\sqrt 6 \qquad \text{(E) }3$

 \textbf{Answer: }C



\textbf{Solution 1}

Each of the numbers $a$ and $b$ is a solution to $\left| x - \frac 1x \right| = 1$.

 Hence it is either a solution to $x - \frac 1x = 1$, or to $\frac 1x - x = 1$. Then it must be a solution either to $x^2 - x - 1 = 0$, or to $x^2 + x - 1 = 0$.

 There are in total four such values of $x$, namely $\frac{\pm 1 \pm \sqrt 5}2$. 

 Out of these, two are positive: $\frac{-1+\sqrt 5}2$ and $\frac{1+\sqrt 5}2$. We can easily check that both of them indeed have the required property, and their sum is $\frac{-1+\sqrt 5}2 + \frac{1+\sqrt 5}2 = \boxed{(C) \sqrt 5}$.

 



\textbf{Solution 2 (speed guess/intuition)}

Since the problem is about similarities a number and its reciprocal differing by one, you can guess that the solution has something to do with the golden ratio, $\varphi = \frac{1+\sqrt 5}2$. Only one of the options has a $\sqrt{5}$ in it, which is $\boxed{(C) \sqrt 5}$.

 


	\item (\textit{AMC 12 12A 2002}) For all positive integers $n$, let $f(n)=\log_{2002} n^2$. Let $N=f(11)+f(13)+f(14)$. Which of the following relations is true?


 $\text{(A) }N<1 \qquad \text{(B) }N=1 \qquad \text{(C) }1<N<2 \qquad \text{(D) }N=2 \qquad \text{(E) }N>2$

 \textbf{Answer: }D



\textbf{Solution}

First, note that $2002 = 11 \cdot 13 \cdot 14$.

 Using the fact that for any base we have $\log a + \log b = \log ab$, we get that $N = \log_{2002} (11^2 \cdot 13^2 \cdot 14^2) = \log_{2002} 2002^2 = \boxed{(D) N=2}$.

 


	\item (\textit{AMC 12 12A 2002}) The mean, median, unique mode, and range of a collection of eight integers are all equal to 8. The largest integer that can be an element of this collection is 


 $\text{(A) }11 \qquad \text{(B) }12 \qquad \text{(C) }13 \qquad \text{(D) }14 \qquad \text{(E) }15$

 \textbf{Answer: }D



\textbf{Solution 1}

As the unique mode is $8$, there are at least two $8$s.

 As the range is $8$ and one of the numbers is $8$, the largest one can be at most $16$.

 If the largest one is $16$, then the smallest one is $8$, and thus the mean is strictly larger than $8$, which is a contradiction.

 If we have 2 8's we can find the numbers 4, 6, 7, 8, 8, 9, 10, 12.This is a possible solution but has not reached the maximum.

 If we have 4 8's we can find the numbers 6, 6, 6, 8, 8, 8, 8, 14.

 We can also see that they satisfy the need for the mode, median, and range to be 8. This means that the answer will be$\boxed{\text{(D)}\ 14 }$. ~By QWERTYUIOPASDFGHJKLZXCVBNM

 



\textbf{Solution 2}

We could express this collection as integers $a\textsubscript{1}$ through $a\textsubscript{8}$, with $a\textsubscript{1}$ being the smallest and $a\textsubscript{8}$ being the largest.

 Since the mode is $8$, we know that $a\textsubscript{4}$ and $a\textsubscript{5}$ must also be $8$. If they were not, the other numbers, which are lesser and greater than $a\textsubscript{4}$ and $a\textsubscript{5}$ respectively, would not be able to satisfy the condition that $8$ is the mode.

 There are $8$ terms and the mean is $8$. This tells us that the sum of all the numbers is $64$.

 We want to maximize the value of $a\textsubscript{8}$, so we set  $a\textsubscript{6}$ and $a\textsubscript{7}$ to $8$ as well.

 Knowing that we want to minimize numbers and that the range is $8$ , we set $a\textsubscript{1}$, $a\textsubscript{2}$, and $a\textsubscript{3}$ equal to $a\textsubscript{8} - 8$.

 {$a\textsubscript{1}$, $a\textsubscript{2}$, $a\textsubscript{3}$, $a\textsubscript{4}$, $a\textsubscript{5}$, $a\textsubscript{6}$, $a\textsubscript{7}$, $a\textsubscript{8}$} $=$ {$a\textsubscript{8} - 8$, $a\textsubscript{8} - 8$, $a\textsubscript{8} - 8$, $8$, $8$, $8$, $8$, $a\textsubscript{8}$}

 Letting the sum of all the numbers be $64$, we find that $32 + 4a_8 - 24 = 64$, which simplifies to $4a_8 = 56$. Solving, we get $\boxed{\text{(D)}\ 14 }$. ~By SK80, mod\_x and huniiya for minor edits

 



\textbf{Solution 3}

Let's say that we choose $k$ numbers to be 8, and the largest number to be $L$. There will therefore be $7-k$ numbers with the value of $L - 8$.

 This gives the following expression $\frac{(7-k)(L-8)+ 8k + L}{8} = 8$. Solving for $k$ gives $\frac{120-8L}{16-L} = k$. The bank of choices we have for $L$ is represented by answer choices (A), (B), (C), (D), and (E).

 We find that the value for $14$ gives $k = 4$ creating the set $666-8888-14$, which satisfies all conditions.

 This gives $\boxed{\text{(D)}\ 14 }$.

 



\textbf{Solution 4}

Since the mean of the numbers is $8$ and there are $8$ numbers, the sum of all $8$ numbers is $64$. Since the range is also $8$, let the smallest value in the set $= x$. Therefore the largest number in the set is $x+8$. The mode is the number which appears most often, so we can let the $4$th to $7$th largest integers $= 8$. To maximize the largest value, we will let the $2$nd and $3$rd smallest integers also equal $x$.

 In order this will be: $x, x, x, 8, 8, 8, 8, x+8$

 Now we can set the expression equal to $64$ (because of what we determined in the $1$st step): $x+x+x+8+8+8+8+x+8=64$

 Simplyifying, 

 $4x+40=64$
$4x=24$
$x=6$

 Therefore, since $x$ is the smallest number in the set, the largest number in the set will be $x+8$, which is $\boxed{\text{(D)}\ 14 }$.

 


	\item (\textit{AMC 12 12A 2002}) Tina randomly selects two distinct numbers from the set $\{1, 2, 3, 4, 5\}$, and Sergio randomly selects a number from the set $\{1, 2, \ldots, 10\}$. What is the probability that Sergio's number is larger than the sum of the two numbers chosen by Tina? 


 $\text{(A)}\ 2/5 \qquad \text{(B)}\ 9/20 \qquad \text{(C)}\ 1/2 \qquad \text{(D)}\ 11/20 \qquad \text{(E)}\ 24/25$

 \textbf{Answer: }A



\textbf{Solution 1}

This is not too bad using casework. 

 Tina gets a sum of 3: This happens in only one way $(1,2)$ and Sergio can choose a number from 4 to 10, inclusive. There are 7 ways that Sergio gets a desirable number here.

 Tina gets a sum of 4: This once again happens in only one way $(1,3)$. Sergio can choose a number from 5 to 10, so 6 ways here.

 Tina gets a sum of 5: This can happen in two ways $(1,4)$ and $(2,3)$. Sergio can choose a number from 6 to 10, so $2\cdot5=10$ ways here.

 Tina gets a sum of 6: Two ways here $(1,5)$ and $(2,4)$. Sergio can choose a number from 7 to 10, so $2\cdot4=8$ here.

 Tina gets a sum of 7: Two ways here $(2,5)$ and $(3,4)$. Sergio can choose from 8 to 10, so $2\cdot3=6$ ways here.

 Tina gets a sum of 8: Only one way possible $(3,5$). Sergio chooses 9 or 10, so 2 ways here.

 Tina gets a sum of 9: Only one way $(4,5)$. Sergio must choose 10, so 1 way.

 In all, there are $7+6+10+8+6+2+1=40$ ways. Tina chooses two distinct numbers in $\binom{5}{2}=10$ ways while Sergio chooses a number in $10$ ways, so there are $10\cdot 10=100$ ways in all. Since $\frac{40}{100}=\frac{2}{5}$, our answer is $\boxed{\text{(A)}\frac{2}{5}}$.

 



\textbf{Solution 2}

We invoke some symmetry. Let $T$ denote Tina's sum, and let $S$ denote Sergio's number. Observe that, for $i = 2, 3, \ldots, 10$, $\text{Pr}(T=i) = \text{Pr}(T=12-i)$.

 If Tina's sum is $i$, then the probability that Sergio's number is larger than Tina's sum is $\frac{10-i}{10}$. Thus, the probability $P$ is

 
 \[P = \text{Pr}(S>T) = \sum_{i=2}^{10} \text{Pr}(T=i) \times \frac{10-i}{10}\]

 Using the symmetry observation, we can also write the above sum as
 \[P = \sum_{i=2}^{10} \text{Pr}(T=12-i) \times \frac{10-i}{10} = \sum_{i=2}^{10} \text{Pr}(T=i) \times \frac{i-2}{10}\]where the last equality follows as we reversed the indices of the sum (by replacing $12-i$ with $i$). Thus, adding the two equivalent expressions for $P$, we have

 
 \begin{align*} 2P &= \sum_{i=2}^{10} \text{Pr}(T=i) \times \left(\frac{10-i}{10} + \frac{i-2}{10}\right) \\ &= \sum_{i=2}^{10} \text{Pr}(T=i) \times \frac{4}{5} \\ &= \frac{4}{5} \sum_{i=2}^{10} \text{Pr}(T=i) \\ &= \frac{4}{5} \end{align*}


 Since this represents twice the desired probability, the answer is $P = \boxed{\textbf{(A)} \frac{2}{5}}$. -scrabbler94

 



\textbf{Solution 3}

We have 4 cases, if Tina chooses $1, 2, 3,$ or  $4$ and always chooses numbers greater than the first number she chose.

 The number of ways of choosing 2 numbers from $5$ are $\binom{5}{2}$.

 Case 1: Tina chooses $1$.

 In this case, since the numbers are distinct, Tina can choose $(1, 2), (1, 3), (1, 4),$ or $(1, 5).$

 If Tina chooses $1$ and $2$ which sum to $3$, Sergio only has $10-3=7$ choices.

 Since the sum of the combined numbers increases by $1$ every time for this specific case, Sergio has $1$ less choice every time.

 Therefore, the probability of this is $\frac{7+6+5+4}{10 \cdot \binom{5}{2}}$.

 Case 2: Tina chooses $2$.

 In this case, Tina can choose $(2, 3), (2, 4),$ or $(2, 5).$

 If Tina chooses $2$ and $3$ which sum to $5$, Sergio only has $10-5=5$ choices.

 Since the sum of the combined numbers increases by $1$ every time for this specific case, Sergio has $1$ less choice every time.

 Therefore, the probability of this is $\frac{5+4+3}{10 \cdot \binom{5}{2}}$.

 Case 3: Tina chooses $3$.

 In this case, Tina can choose $(3, 4),$ or $(3, 5).$

 If Tina chooses $3$ and $4$ which sum to $7$, Sergio only has $10-7=3$ choices.

 Since the sum of the combined numbers increases by $1$ every time for this specific case, Sergio has $1$ less choice every time.

 Therefore, the probability of this is $\frac{3+2}{10 \cdot \binom{5}{2}}$.

 Case 4: Tina chooses $4$.

 In this case, Tina can only choose $(4,5).$

 If Tina chooses $4$ and $5$ which sum to $9$, Sergio only has $10-9=1$ choice.

 Therefore, the probability of this is $\frac{1}{10 \cdot \binom{5}{2}}$.

 Once you add these probabilities up, you will have $\frac{(7+6+5+4)+(5+4+3)+(3+2)+(1)}{10 \cdot\binom{5}{2}} = \frac{40}{100} = \frac{2}{5}$ probability.

 Thus our answer is $\frac{2}{5}$.

 



\textbf{Solution 4}

Assume Sergio chooses from ${2,3,\ldots,10}$. The probability of Tina getting a sum of $6+x$ and a sum of $6-x$, where $x \leq 4$, are equal due to symmetry. The probability of Sergio choosing numbers higher/lower than $6+x$ is equal to him choosing numbers lower/higher than $6-x$. Therefore over all of Tina's sums, the probability of Sergio choosing a number higher is equal to the probability of choosing a number lower. 

 The probability that they get the same value is $1/9$, so the probability of Sergio getting a higher number is $\frac{(9-1)/2}{9} = \frac49$. 

 Sergio never wins when choosing $1$ so the probability is $\frac49 \cdot \frac{9}{10} + (0)\frac{1}{10} = \boxed{\textbf{(A)} \frac{2}{5}}.$




\textbf{Solution 5 (Brute Force)}

List all the cases where $S \in [1, 10]$ and you get $\frac{0+0+0+1+2+4+6+8+9+10}{\binom{5}{2} \cdot 10} = \boxed{\textbf{(A)} \frac{2}{5}}$




\textbf{Solution 6 (Cheese)}

Add up the first set and then divide by $5$ and you get $3$. Multiply this by two and you can see the expected value of the sum is $6$. There is a $\frac{2}{5}$ chance you will get a number higher this is great if you are running low on time.

 (Note: this method only works because Tina's largest value ($9$) is less than Sergio's ($10$). If Sergio chose a random number from a set like $\{1, 2, 3, 4, 5, 6, 7\}$, this method gives an answer of $14.3 \%$ whereas the actual answer would be $18.6 \%$.)

 


	\item (\textit{AMC 12 12A 2002}) Several sets of prime numbers, such as $\{7,83,421,659\}$ use each of the nine nonzero digits exactly once. What is the smallest possible sum such a set of primes could have?


 $\text{(A) }193 \qquad \text{(B) }207 \qquad \text{(C) }225 \qquad \text{(D) }252 \qquad \text{(E) }447$

 \textbf{Answer: }B



\textbf{Solution}

Neither of the digits $4$, $6$, and $8$ can be a units digit of a prime. Therefore the sum of the set is at least $40 + 60 + 80 + 1 + 2 + 3 + 5 + 7 + 9 = 207$.

 We can indeed create a set of primes with this sum, for example the following sets work: $\{ 41, 67, 89, 2, 3, 5 \}$ or $\{ 43, 61, 89, 2, 5, 7 \}$.

 Thus the answer is $207\implies \boxed{\mathrm{(B)}}$.

 


	\item (\textit{AMC 12 12A 2002}) Let $C_1$ and $C_2$ be circles defined by $(x-10)^2 + y^2 = 36$ and $(x+15)^2 + y^2 = 81$
respectively. What is the length of the shortest line segment $PQ$ that is tangent to $C_1$ at $P$ and to $C_2$ at $Q$?


 $\text{(A) }15 \qquad \text{(B) }18 \qquad \text{(C) }20 \qquad \text{(D) }21 \qquad \text{(E) }24$

 \textbf{Answer: }C



\textbf{Solution 1}

First examine the formula $(x-10)^2+y^2=36$, for the circle $C_1$. Its center, $D_1$, is located at (10,0) and it has a radius of $\sqrt{36}$ = 6. The next circle, using the same pattern, has its center, $D_2$, at (-15,0) and has a radius of $\sqrt{81}$ = 9. So we can construct this diagram:
                    \begin{figure}[H]
                        \centering
                        \includegraphics[width=0.5\textwidth]{Solution/AMC_12_12A_2002_18_Solution_Figure_1.png}
                    \end{figure}
                    Line PQ is tangent to both circles, so it forms a right angle with the radii (6 and 9). This, as well as the two vertical angles near O, prove triangles $D_2QO$ and $D_1PO$ similar by AA, with a scale factor of $6:9$, or $2:3$. Next, we must subdivide the line $D_2D_1$ in a 2:3 ratio to get the length of the segments $D_2O$ and $D_1O$. The total length is $10 - (-15)$, or $25$, so applying the ratio, $D_2O$ = \textbf{15} and $D_1O$ = \textbf{10}. These are the hypotenuses of the triangles. We already know the length of $D_2Q$ and $D_1P$, \textbf{9} and \textbf{6} (they're radii). So in order to find $PQ$, we must find the length of the longer legs of the two triangles and add them. 

 $15^2 - 9^2 = (15-9)(15+9) = 6 \times 24 = 144$
$\sqrt{144} = 12$
$10^2-6^2 = (10-6)(10+6) = 4 \times 16 = 64$
$\sqrt{64} = 8$

 Finally, the length of PQ is $12+8=\boxed{20}$, or \textbf{(C)}.

 



\textbf{Solution 2}

Using the above diagram, imagine that segment $\overline{QD_2}$ is shifted to the right to match up with $\overline{PD_1}$. Then shift $\overline{QP}$ downwards to make a right triangle. We know $\overline{D_2D_1} = 25$ from the given information and the newly created leg has length $\overline{QD_2} + \overline{PD_1} = 9 + 6 = 15$. Hence by Pythagorean theorem $15^2 + {\overline{QP}}^2 = 25^2$.

 $\overline{QP} = \boxed{20}$, or C.

 


	\item (\textit{AMC 12 12A 2002}) The graph of the function $f$ is shown below. How many solutions does the equation $f(f(x))=6$ have? 


 
                    \begin{figure}[H]
                        \centering
                        \includegraphics[width=0.5\textwidth]{Problem/AMC_12_12A_2002_19_Problem_Figure_1.png}
                    \end{figure}
                    
$\text{(A) }2 \qquad \text{(B) }4 \qquad \text{(C) }5 \qquad \text{(D) }6 \qquad \text{(E) }7$

 \textbf{Answer: }D



\textbf{Solution}

First of all, note that the equation $f(t)=6$ has two solutions: $t=-2$ and $t=1$. 

 Given an $x$, let $f(x)=t$. Obviously, to have $f(f(x))=6$, we need to have $f(t)=6$, and we already know when that happens. In other words, the solutions to $f(f(x))=6$ are precisely the solutions to ($f(x)=-2$ or $f(x)=1$).

 Without actually computing the exact values, it is obvious from the graph that the equation $f(x)=-2$ has two and $f(x)=1$ has four different solutions, giving us a total of $2+4=\boxed{(D)6}$ solutions.

 
                    \begin{figure}[H]
                        \centering
                        \includegraphics[width=0.5\textwidth]{Solution/AMC_12_12A_2002_19_Solution_Figure_1.png}
                    \end{figure}
                    



	\item (\textit{AMC 12 12A 2002}) Suppose that $a$ and $b$ are digits, not both nine and not both zero, and the repeating decimal $0.\overline{ab}$ is expressed as a fraction in lowest terms. How many different denominators are possible?


 $\text{(A) }3 \qquad \text{(B) }4 \qquad \text{(C) }5 \qquad \text{(D) }8 \qquad \text{(E) }9$

 \textbf{Answer: }C



\textbf{Solution 1}

The repeating decimal $0.\overline{ab}$ is equal to 
 \[\frac{10a+b}{100} + \frac{10a+b}{10000} + \cdots = (10a+b)\cdot\left(\frac 1{10^2} + \frac 1{10^4} + \cdots \right) =  (10a+b) \cdot \frac 1{99} = \frac{10a+b}{99}\]

 When expressed in the lowest terms, the denominator of this fraction will always be a divisor of the number $99 = 3\cdot 3\cdot 11$. This gives us the possibilities $\{1,3,9,11,33,99\}$. As $a$ and $b$ are not both nine and not both zero, the denominator $1$ can not be achieved, leaving us with $\boxed{\mathrm{(C) }5}$ possible denominators.

 (The other ones are achieved e.g. for $ab$ equal to $33$, $11$, $9$, $3$, and $1$, respectively.)

 



\textbf{Solution 2}

Another way to convert the decimal into a fraction (simplifying, I guess?). We have 
 \[100(0.\overline{ab}) = ab.\overline{ab}\] 
 \[99(0.\overline{ab}) = 100(0.\overline{ab}) - 0.\overline{ab} = ab.\overline{ab} - 0.\overline{ab} = ab\] 
 \[0.\overline{ab} = \frac{ab}{99}\]where $a, b$ are digits. Continuing in the same way by looking at the factors of 99, we have 5 different possibilities for the denominator. $\boxed{(C)}$




\textbf{Solution 3}

Since $\frac{1}{99}=0.\overline{01}$, we know that $0.\overline{ab} = \frac{ab}{99}$. From here, we wish to find the number of factors of $99$, which is $6$. However, notice that $1$ is not a possible denominator, so our answer is $6-1=\boxed{5}$.
 \[\]~AopsUser101

 



\textbf{Solution 4 (Alcumus)}

Since $0.\overline{ab} = \frac{ab}{99}$, the denominator must be a factor of $99 = 3^2 \cdot 11$. The factors of $99$ are $1,$ $3,$ $9,$ $11,$ $33,$ and $99$. Since $a$ and $b$ are not both nine, the denominator cannot be $1$. By choosing $a$ and $b$ appropriately, we can make fractions with each of the other denominators.

 Thus, the answer is $\boxed{5}$.

 


	\item (\textit{AMC 12 12A 2002}) Consider the sequence of numbers: $4,7,1,8,9,7,6,\dots$ For $n>2$, the $n$-th term of the sequence is the units digit of the sum of the two previous terms. Let $S_n$ denote the sum of the first $n$ terms of this sequence. The smallest value of $n$ for which $S_n>10,000$ is: 


 $\text{(A) }1992 \qquad \text{(B) }1999 \qquad \text{(C) }2001 \qquad \text{(D) }2002 \qquad \text{(E) }2004$

 \textbf{Answer: }B



\textbf{Solution}

The sequence is infinite. As there are only $100$ pairs of digits, sooner or later a pair of consecutive digits will occur for the second time. As each next digit only depends on the previous two, from this point on the sequence will be periodic.

 (Additionally, as every two consecutive digits uniquely determine the \textit{previous} one as well, the first pair of digits that will occur twice must be the first pair $4,7$.)

 Hence it is a good idea to find the period. Writing down more terms of the sequence, we get:

 
 \[4,7,1,8,9,7,6,3,9,2,1,3,4,7,\dots\]

 and we found the period. The length of the period is $12$, and its sum is $4+7+\cdots+1+3 = 60$. Hence for each $k$ we have $S_{12k} = 60k$.

 We have $\lfloor 10000/60 \rfloor = 166$ and $166\cdot 12 = 1992$, therefore $S_{1992} = 60\cdot 166 = 9960$.The rest can now be computed by hand, we get $S_{1998} = 9960+4+7+1+8+9+7= 9996$, and $S_{1999}=9996 + 6 = 10002$, thus the answer is $\boxed{\text{(B) }1999}$.

 


	\item (\textit{AMC 12 12A 2002}) Triangle $ABC$ is a right triangle with $\angle ACB$ as its right angle, $m\angle ABC = 60^{\circ}$, and $AB = 10$. Let $P$ be randomly chosen inside $\triangle ABC$, and extend $\overline{BP}$ to meet $\overline{AC}$ at $D$. What is the probability that $BD > 5\sqrt{2}$?


 $\textbf{(A)}\ \frac{2-\sqrt2}{2}\qquad\textbf{(B)}\ \frac{1}{3}\qquad\textbf{(C)}\ \frac{3-\sqrt3}{3}\qquad\textbf{(D)}\ \frac{1}{2}\qquad\textbf{(E)}\ \frac{5-\sqrt5}{5}$

 \textbf{Answer: }C



\textbf{Solution}

Clearly $BC=5$ and $AC=5\sqrt{3}$. Choose a $P'$ and get a corresponding $D'$ such that $BD'= 5\sqrt{2}$ and $CD'=5$. For $BD > 5\sqrt2$ we need $CD>5$, creating an isosceles right triangle with hypotenuse $5\sqrt {2}$ . Thus the point $P$ may only lie in the triangle $ABD'$. The probability of it doing so is the ratio of areas of $ABD'$ to $ABC$, or equivalently, the ratio of $AD'$ to $AC$ because the triangles have identical altitudes when taking $AD'$ and $AC$ as bases. This ratio is equal to $\frac{AC-CD'}{AC}=1-\frac{CD'}{AC}=1-\frac{5}{5\sqrt{3}}=1-\frac{\sqrt{3}}{3}= \frac{3-\sqrt{3}}{3}$. Thus the answer is $\boxed{C}$.

 


	\item (\textit{AMC 12 12A 2002}) In triangle $ABC$, side $AC$ and the perpendicular bisector of $BC$ meet in point $D$, and $BD$ bisects $\angle ABC$.  If $AD = 9$ and $DC = 7$, what is the area of triangle $ABD$?


 $\text{(A)}\ 14 \qquad \text{(B)}\ 21 \qquad \text{(C)}\ 28 \qquad \text{(D)}\ 14\sqrt5 \qquad \text{(E)}\ 28\sqrt5$

 \textbf{Answer: }D



\textbf{Solution 1}


                    \begin{figure}[H]
                        \centering
                        \includegraphics[width=0.5\textwidth]{Solution/AMC_12_12A_2002_23_Solution_Figure_1.png}
                    \end{figure}
                    Looking at the triangle $BCD$, we see that its perpendicular bisector reaches the vertex, therefore implying it is isosceles. Let $x = \angle C$, so that $B=2x$ from given and the previous deducted. Then $\angle ABD=x, \angle ADB=2x$ because any exterior angle of a triangle has a measure that is the sum of the two interior angles that are not adjacent to the exterior angle. That means $\triangle ABD$ and $\triangle ACB$ are \href{https://artofproblemsolving.com/wiki/index.php?title=Similarity\_(geometry)}{similar}, so $\frac {16}{AB}=\frac {AB}{9} \Longrightarrow AB=12$.

 Then by using \href{https://artofproblemsolving.com/wiki/index.php?title=Heron\%27s\_Formula}{Heron's Formula} on $ABD$ (with sides $12,7,9$), we have $[\triangle ABD]= \sqrt{14(2)(7)(5)} = 14\sqrt5 \Longrightarrow \boxed{\text{D}}$.

 



\textbf{Solution 2}

Let M be the point of the perpendicular bisector on BC. By the perpendicular bisector theorem, $BD = DC = 7$ and $BM = MC$. Also, by the angle bisector theorem, $\frac {AB}{BC} = \frac{9}{7}$. Thus, let $AB = 9x$ and $BC = 7x$. In addition, $BM = 3.5x$. 

 Thus, $\cos\angle CBD = \frac {3.5x}{7} = \frac {x}{2}$. Additionally, using the Law of Cosines and the fact that $\angle CBD = \angle ABD$, $81 = 49 + 81x^2 - 2(9x)(7)\cos\angle CBD$

 Substituting and simplifying, we get $x = 4/3$

 Thus, $AB = 12$. We now know all sides of $\triangle ABD$. Using \href{https://artofproblemsolving.com/wiki/index.php?title=Heron\%27s\_Formula}{Heron's Formula} on $\triangle ABD$, $\sqrt{(14)(2)(7)(5)} = 14\sqrt5 \Longrightarrow \boxed{\text{D}}$

 "Note:-you could also drop a perpendicular from D to AB at point let say,F then BF = 3.5x by pyathgoras theorem we can find DF and (AB ×DF )÷2 is our answer"

 



\textbf{Solution 3}

Note that because the perpendicular bisector and angle bisector meet at side $AC$ and $CD = BD$ as triangle $BDC$ is isosceles, so $BD = 7$. By the angle bisector theorem, we can express $AB$ and $BC$ as $9x$ and $7x$ respectively. We try to find $x$ through Stewart's Theorem. So

 $16(7^2+9\cdot7) = (7x)^2 \cdot 9 + (9x)^2 \cdot 7$
$16(49+63) = (49 \cdot 9 + 81 \cdot 7)x^2$
$16(49+63) = 9(49+63) \cdot x^2$
$x^2 = \frac{16}{9}$
$x=\frac{4}{3}$

 We plug this to find that the sides of $\triangle ABD$ are $12,7,9$. By Heron's formula, the area is  $\sqrt{(14)(2)(7)(5)} = 14\sqrt5 \Longrightarrow \boxed{\text{D}}$. ~skyscraper

 



\textbf{Solution 4}


                    \begin{figure}[H]
                        \centering
                        \includegraphics[width=0.5\textwidth]{Solution/AMC_12_12A_2002_23_Solution_Figure_2.png}
                    \end{figure}
                    

 Draw $DN$ such that $DN \bot AB$, $\triangle BND \cong \triangle BMD$
$\angle ACB = \angle DBC = \angle ABD$, $\triangle ABD \sim \triangle ACB$ by $AA$
$\frac{AB}{AD} = \frac{AC}{AB}$, $AB^2 = AD \cdot AC = 9 \cdot 16$, $AB = 12$

 By the \href{https://artofproblemsolving.com/wiki/index.php?title=Angle\_Bisector\_Theorem}{Angle Bisector Theorem}, $\frac{BC}{AB} = \frac{CD}{AD}$
$\frac{BC}{12} = \frac{7}{9}$
$BC = \frac{28}{3}$, $CM = \frac{14}{3}$, $DN = DM = \sqrt{CD^2 - CM^2} = \frac{7 \sqrt{5} }{3}$
$[ABD] = \frac12 \cdot AB \cdot DN = \frac12 \cdot 12 \cdot \frac{7 \sqrt{5} }{3} = \boxed{\textbf{(D) } 14 \sqrt{5} }$




\textbf{Solution 5 (Trigonometry)}

Let $\angle ACB = \theta$, $\angle DBC = \theta$, $\angle ABD = \theta$, $\angle ADB = 2 \theta$, $\angle BAC = 180^\circ - 3 \theta$
$[ABD] = \frac12 \cdot AD \cdot BD \cdot \sin \angle ADB = \frac12 \cdot 9 \cdot 7 \cdot \sin 2\theta$

 By the \href{https://artofproblemsolving.com/wiki/index.php?title=Law\_of\_Sines}{Law of Sines} we have $\frac{ \sin \angle BAC }{BD} = \frac{ \sin \angle ABD }{AD}$
$\frac{ \sin(180^\circ - 3 \theta) }{7} = \frac{ \sin \theta }{ 9 }$
$\frac{ \sin 3 \theta }{7} = \frac{ \sin \theta }{ 9 }$

 By the \href{https://artofproblemsolving.com/wiki/index.php/Trigonometric\_identities\#Triple-angle\_identities}{Triple-angle Identities}, $\sin 3 \theta = 3 \sin \theta - 4 \sin^3 \theta$
$3 - 4 \sin^2 \theta = \frac79$, $36 \sin^2 \theta = 20$, $\sin^2 \theta = \frac59$
$\sin \theta = \frac{\sqrt{5}}{3}$, $\cos \theta = \frac{2}{3}$

 By the \href{https://artofproblemsolving.com/wiki/index.php/Trigonometric\_identities\#Double-angle\_identities}{Double Angle Identity} $\sin 2 \theta = 2 \sin \theta \cos \theta = 2 \cdot \frac{\sqrt{5}}{3} \cdot \frac{2}{3} = \frac{ 4 \sqrt{5} }{9}$
$[ABD] = \frac12 \cdot 9 \cdot 7 \cdot \frac{ 4 \sqrt{5} }{9} =  \boxed{\textbf{(D) } 14 \sqrt{5} }$



	\item (\textit{AMC 12 12A 2002}) Find the number of ordered pairs of real numbers $(a,b)$ such that $(a+bi)^{2002} = a-bi$.


 $\text{(A) }1001 \qquad \text{(B) }1002 \qquad \text{(C) }2001 \qquad \text{(D) }2002 \qquad \text{(E) }2004$

 \textbf{Answer: }E



\textbf{Solution 1}

Let $s=\sqrt{a^2+b^2}$ be the magnitude of $a+bi$. Then the magnitude of $(a+bi)^{2002}$ is $s^{2002}$, while the magnitude of $a-bi$ is $s$. We get that $s^{2002}=s$, hence either $s=0$ or $s=1$.

 For $s=0$ we get a single solution $(a,b)=(0,0)$.

 Let's now assume that $s=1$. Multiply both sides by $a+bi$. The left hand side becomes $(a+bi)^{2003}$, the right hand side becomes $(a-bi)(a+bi)=a^2 + b^2 = 1$. Hence the solutions for this case are precisely all the $2003$rd complex roots of unity, and there are $2003$ of those.

 The total number of solutions is therefore $1+2003 = \boxed{2004}$.

 



\textbf{Solution 2}

As in the other solution, split the problem into when $s=0$ and when $s=1$. When $s=1$ and $a+bi=\cos\theta+i\sin\theta$,

 $(a+bi)^{2002}= \cos(2002\theta)+i\sin(2002\theta)$ $=a-bi= \cos\theta-i\sin\theta= \cos(-\theta)+i\sin(-\theta)$

 so we must have $2002\theta=-\theta+2\pi k$ and hence $\theta=\frac{2\pi k}{2003}$. Since $\theta$ is restricted to $[0,2\pi)$, $k$ can range from $0$ to $2002$ inclusive, which is $2002-0+1=2003$ values. Thus the total is $1+2003 = \boxed{\textbf{(E)}\  2004}$.

 



\textbf{Solution 3}

Let $r$ be the magnitude of $a+bi$. Notice that $r$ must be either $0$ or $1$ for this to be true, as shown in the above solutions. We know this because we are taking magnitude to the $2003$rd power, and if the magnitude of $a+bi$ is larger than $1$, it will increase and if it is smaller than $1$ it will decrease. However, the magnitude on the RHS is still $r$, so this is not possible. Again, only $r=0$ and $r=1$ satisfy the equation.

 Now if $r=0$, then $(a,b)$ must be $(0,0)$.

 However if $r=1$, we then have:

 $\cos(2002 \theta) = \cos(-\theta)$. This has solution of $\theta = 0$. This would represent the number $1+0i$, with conjugate $1-0i$. This works because the magnitude is the same and the angle is nothing anyways. We multiply angle by $2002$ through De Moivre's Theorem and also we do $-\theta$ because it is a reflection, angles therefore is negative.

 We then write:

 $\cos(2002 \theta) = \cos(360-\theta)$ which has solution of $\theta = \frac{360}{2003}$. 

 We can also write:

 $\cos(2002 \theta) = \cos(720-\theta)$ which has solution $\theta = \frac{720}{2003}$.

 We notice that it is simply headed upwards and the answer is of the form $\frac{720}{2003} n$, where n is some integer from $0$ to infinity, inclusive.

 Well wait, it repeats itself $n=2003$, that is $360$ which is also $0$! Hence we only have $n=0$ to $2002$ as original solutions, or $2003$ solutions.

 $1+2003 = \boxed{2004}$.

 
Solution by Blackhawk 9-10-17

 $\LaTeX$ed (with some edits) by PhunsukhWangdu 7/27/22

 


 



\textbf{Solution 4}

Let $z = a + bi = re^{i\theta}$ where $r = |z|$ and $0\leq \theta < 2\pi.$ We want to solve $z^{2002} = \overline{z}.$ Since $z\overline{z} = |z|^2,$ we multiply by $z$ on both sides to get 
 \[z^{2003} = |z|^2 = r^{2003}e^{2023i\theta} = r^2,\]from which we get $r\in \{0, 1\}$ (since $r\in \mathbb{R}$), immediately yielding $z = 0$ as an answer. If $r=1$ then we have $z^{2003} = 1$, which each of the $2003$rd roots of unity satisfies. Altogether, there are $1 + 2003 = \boxed{\textbf{(E)\ 2004}}$ values of $z$, each of which correspond to a unique ordered pair $(a, b)$. 

 -Benedict T(countmath1)

 



\textbf{Solution 5}

Let $z=a+bi \Rightarrow z^{2022}=\overline{z}.$ Hence, taking the magnitude of both sides, we must have:
 \[|z^{2022}|=|\overline{z}|=|z| \Rightarrow |z|=0, 1.\]If $|z|=0,$ then $z=0.$ This gives one solution.

 If $|z|=1,$ then we do the following. Multiply the first equation by $z$ on both sides such that you obtain $z^{2003}=z \overline{z} =|z|^2.$ 

 Now, it is clear that if $|z|=1,$ we need $z^{2003}=1.$ We have 2003 solutions for this, i.e. for each $i\in\{0,1,\cdots,2022\}$ for $e^{2i\pi/2003}.$

 A total of $2003+1=2004$ solutions.

 


	\item (\textit{AMC 12 12A 2002}) The nonzero coefficients of a polynomial $P$ with real coefficients are all replaced by their mean to form a polynomial $Q$. Which of the following could be a graph of $y = P(x)$ and $y = Q(x)$ over the interval $-4\leq x \leq 4$?


 
                    \begin{figure}[H]
                        \centering
                        \includegraphics[width=0.5\textwidth]{Problem/AMC_12_12A_2002_25_Problem_Figure_1.png}
                    \end{figure}
                    

 \textbf{Answer: }B



\textbf{Solution 1}

The sum of the coefficients of $P$ and of $Q$ will be equal, so $P(1) = Q(1)$. The only answer choice with an intersection between the two graphs at $x = 1$ is \textbf{(B)}. (The polynomials in the graph are $P(x) = 2x^4-3x^2-3x-4$ and $Q(x) = -2x^4-2x^2-2x-2$.)

 



\textbf{Solution 2}

We know every coefficient is equal, so we get $ax^n + ... + ax + a = 0$ which equals $x^n + ... + x + 1 = 0$. We see apparently that x cannot be positive, for it would yield a number greater than zero for $Q(x)$. We look at the zeros of the answer choices. A, C, D, and E have a positive zero, which eliminates them. B is the answer.

 



\textbf{Solution 2a}

Proceed similarly to Solution 2. However, note that $x^n + ... + x + 1$ is a factor of $x^{n+1}-1.$ The only potential zeroes of $x^{n+1}-1$ are $1$ and $-1$ if $n+1$ is even.  Since none of the graphs have a zero at exactly negative one, the solution is the graph with no zeroes, which is $\boxed{(B)}.$



\end{enumerate}